% ===== §27 (new: family ethos bridge case) =====
\section{Family Ethos as a Middle-Scale Case: \emph{Jiafeng}}
\label{sec:familyethos}

\noindent
The argument has established a general ontology of cultural emergence and a phenomenology of its most intimate instance, and it is about to turn to the political economy of production and capture. Between these stands a case worth examining in its own right, because it occupies a scale the paper has so far left empty: the scale between the intimate dyad and the macro-culture of a people. That case is the \emph{family ethos}---in Chinese \emph{jiafeng} (\zh{家风}), the ``wind'' or manner of a household, the distinctive culture a family generates and transmits across generations. The family ethos is chosen not as an illustration of the theory but as an object through which two questions the theory has raised but not settled can be worked out: how a culture is to be \emph{known} historically, which is a question of the epistemology of a dialectical view; and how power \emph{constitutes} a culture without exhausting it, which is a question of what may be called the study of power in culture. The case is the thread; the epistemology and the power-analysis are drawn from it as they arise.

\subsection*{The family ethos as a middle-scale culture}

The family ethos sits, in the terms of this paper, at a scale between the two the argument has treated. It is more durable and more transmissible than the culture of a dyad, for it passes across generations and outlives the particular members who bear it at any moment; and it is more endogenous than the culture of a people, for it is generated within a bounded relation---the family---that can author its own manner in a way a whole society cannot. It has both faces the paper's opening set in tension: a \emph{wen}, the sedimented pattern of transmissible forms, and a \emph{hua}, the ongoing cultivation by which each generation is formed and forms the next. Its designs, across civilizations, fall into recognizable types. There is the \emph{textual} type, which sediments the ethos into written instruction: the Chinese tradition of \emph{jiaxun} (\zh{家训}), family instructions, of which Yan Zhitui's sixth-century \emph{Family Instructions for the Yan Clan} is the earliest surviving and most influential, a book-length letter to his sons on conduct, learning, and the management of a household \citep{yanzhitui2021}. There is the \emph{unwritten-exemplary} type, which transmits the ethos not by text but by modeled example and remembered saying: the Roman \emph{mos maiorum}, the ``way of the ancestors,'' an unwritten code carried across generations through ritual, oratory, and the exemplary deeds of the \emph{maiores}, in dynamic complement to written law rather than as a codification of it. And there is the \emph{virtue-and-model} type, which sediments the ethos into a named set of dispositions embodied in the elder---diligence, frugality, honesty, the ``farming-and-reading'' ideal of the scholar household---transmitted less by precept than by the standing example of a life. These types are not exclusive, and most family ethical traditions mix them; they are named here to show the range of designs by which a family gives its manner a transmissible form.

\subsection*{The dialectical evaluation, and the epistemology it requires}

The family ethos is the ideal case for the disciplined historical evaluation the synthesis required, because its forms are old enough to have changed their meaning as their conditions changed, and evaluating them tests what a dialectical view of culture actually demands of the one who holds it. Consider the ``farming-and-reading'' ethos, the ideal of the household that works the land and cultivates learning together. Under the material conditions that produced it---a smallholding agrarian economy with a channel of advancement through examination and letters---this ethos was a genuine advance: it encoded a real emancipatory impulse, the aspiration of a family to lift itself, through cultivation, beyond the bare necessity of the soil, and it generated, in those conditions, more than it foreclosed. To evaluate it only by a later standard, as a quaint or constraining traditionalism, would be to commit the error a dialectical view forbids. Consider, on the other side, an ethos organized around a strict hierarchy of elder over younger, of one generation's authority over the next. Under the conditions of a lineage-based agrarian order, such an ethos stabilized the reproduction of the household and had, in its setting, a coordinating function that was not mere domination. Yet as the material conditions change---as economies industrialize, as households become mobile and nuclear, as the young acquire, through education and independent livelihood, a standing the old order did not grant them---the same ethos that once coordinated can become a fetter, its hierarchy losing the material basis that once made it functional and persisting as bare constraint. The one ethos was progressive under its conditions and is to be criticized under changed ones; and to criticize it now and to acknowledge its former progressiveness are, as the framework held, two faces of one evaluation, not a contradiction.

This is the point at which the case yields its first theoretical moment, the \emph{epistemology of a dialectical view}. To evaluate a family ethos as this evaluation requires is to occupy a difficult epistemic position, defined by its refusal of two easier ones. It refuses \emph{presentism}, the judging of a past form by present standards as though its conditions were irrelevant, which would convict every old ethos of failing to be modern. And it refuses the \emph{historicist relativism} that would say a form can be understood only within its own conditions and never evaluated across them, which would leave one unable to criticize any ethos at all and would collapse into the position that whatever was, under its conditions, was right. The dialectical view is the third position, and its possibility rests on a specific ground: that the standard of evaluation is neither an external absolute imported from the present nor the form's own self-image, but \emph{generativity} itself---whether the ethos, under its conditions, let the family generate more than it foreclosed, and whether the power it carried was met by a force that kept it from entrenching. This standard is internal enough to respect the form's conditions and general enough to evaluate across them, and it is what lets the dialectical view avoid both presentism and relativism. The epistemology has a further requirement, which the debate over the fate of the Confucian tradition makes visible. Where one influential reading held the incongruity between Confucian traditionalism and modernity to be so clear-cut that the coming of the modern entailed the tradition's demise \citep{levenson1968}, another held that the tradition does not simply die but undergoes transformation, passing through epochs in which it is reinterpreted rather than abolished \citep{tu1986}. The dialectical view sides with neither the thesis of demise nor the thesis of an unchanging essence; it holds that a cultural form is known rightly only from the tension of two temporal positions at once---the position of its own moment, in which it was the negation its conditions required, and the position of the later moment, in which it becomes the negation to be negated. To know a family ethos historically is to hold both positions together, and the refusal to hold both is the epistemic error that produces, on one side, the verdict of mere demise, and on the other, the illusion of a timeless tradition.

One further epistemic point the case forces into view: the one who evaluates a family ethos is not outside it. The evaluator was formed by some family ethos, shaped in their very capacity to evaluate by a culture of the kind they now assess, and the dialectical view must include this reflexive recognition on pain of relapsing into the idealist fiction of a knower standing outside history. This connects the epistemology of the dialectical view to the account of participatory knowing given earlier: as intimate culture is known only from within its circuit, the historical knowing of a family ethos is carried on by a knower who is themselves a product of such cultures, and whose evaluation is therefore a moment in the very history it evaluates, not a view from beyond it.

\subsection*{Power in culture, and the force that re-orients it}

The case yields its second theoretical moment when the same ethos is examined for how power runs through it, and this is the point at which the family ethos becomes a middle-scale test of the paper's central claim about power and culture. On one side, the family ethos confirms, with unusual clarity, the thesis that power is internal to culture: an ethos of generational hierarchy plainly encodes a power order, naturalizing an authority of elder over younger into the very manner of the household, so that the power tradition's insight is vindicated at this scale as much as at the scale of a people. A family ethos is, among other things, an encoding of who defers to whom. On the other side, the family ethos equally confirms the limit of that thesis, for an ethos is manifestly not reducible to its encoding of power: the same ethos carries real cultivation, genuine virtues, an actual formation of persons, and a real bond of familial love that belongs to the register of the relational Real, and no analysis that dissolved the ethos wholly into a power order would capture what it is. The family ethos thus stands as a middle-scale confirmation of both the power tradition and its sublation: power constitutes the ethos, and does not exhaust it, because the ethos spans registers a pure analysis of power cannot reach.

Between these, the case makes legible---more legible than the dyad could---the distinction between power that entrenches and power that is kept open, and it yields, in doing so, a generalizable proposition about power in culture. A family ethos that encodes only a hierarchy, only the question of who obeys whom, is an ethos of \emph{power-over}: it fixes the young into a determined position and reproduces, generation upon generation, the authority it naturalizes, and such an ethos, absent any internal check, tends toward the pure reproduction of domination. But some family ethical traditions encode, alongside the authority they transmit, a \emph{force that re-orients that authority}: the ``reading'' in the farming-and-reading ideal is exactly such a force, for the cultivation of learning gives the next generation the capacity to question, to exceed, and to re-orient the very ethos that formed them; and the place some traditions reserve for remonstrance---the recognized standing of the younger to admonish the elder---institutionalizes, within the ethos, a counter-power against the entrenchment of the authority the ethos also carries. Such an ethos is one of \emph{power-to}: it transmits authority and, in the same transmission, transmits the means of that authority's own revision. And the historical observation the case supports is precise and generalizable beyond the family: the family ethical traditions that have persisted across changing conditions without hardening into mere oppression are not, on the whole, the most rigid, but those that encoded within themselves a mechanism of their own revision---a force by which each generation could re-orient the ethos to conditions its authors could not foresee. An ethos that transmits only its authority entrenches and, when conditions change, becomes the fetter that must be broken; an ethos that transmits also the power to revise it develops, bending to new conditions without breaking. This is the sublation of power, established in the synthesis for culture in general, confirmed at the scale of the family: the just and the durable ethos is not the one without power, which does not exist, but the one that carries, with its power, the counter-power that keeps that power from entrenching.

From this the general study of power in culture takes a methodological principle, which the case has earned rather than assumed. The task of a critical analysis of culture is not to seek a culture free of power, which is a naivety the power tradition rightly dismisses, nor to conclude, from power's ubiquity, that all culture is equally and irredeemably domination, which is the opposite error. It is to discriminate, within a culture shot through with power as all culture is, between the entrenchment of power and its re-orientation---to ask, of any cultural form, not whether it contains power but whether it contains, with its power, the force that keeps that power open. The family ethos, examined at its own scale, is where this principle becomes visible, because in the family one can see, across generations, both the ethos that hardened into the reproduction of an authority outlived by its conditions, and the ethos that carried within it the means of its own generational revision. The middle scale is the scale at which the difference between power-over and power-to in culture is most clearly read, and reading it there confirms, for culture in general, that the presence of power is not the end of the analysis but its beginning.

\subsection*{The bridge to production and capture}

The family ethos, finally, returns the argument to the scale question and hands it forward to the political economy that follows. It shows, against the weak analogy the paper has held to, that endogenous cultural generation does occur above the scale of the dyad---that a family, and not only a couple, can author a culture of its own, generated within its bounded relation and transmitted across its generations. It does not show that the same holds at the scale of a whole symbolic order, and the caution the paper has reserved for that further step remains in force; the family is a middle term, not a bridge all the way to the macro-scale. But it establishes that the endogenous is not confined to the smallest unit, and that the dialectic of power and counter-power the synthesis described operates wherever culture is generated, at the scale of the family as at the scale of the couple. And it hands the argument forward, for the family ethos is produced under material conditions, as the next section will insist of all culture, and is, in the present, as subject to capture as any other endogenous culture---its manner supplied ready-made by a market in family lifestyles, its cultivation outsourced to purchased instruction, its transmission mediated by platforms that own the means through which the generations now communicate. The production and capture of culture, to which the argument now turns, reach the family ethos as they reach the culture of the dyad, and the case just examined becomes, in the sections that follow, one more site at which the endogenous generation of culture meets the material conditions that constrain it and the capital that would draw its value off.
